This chapter describes the measurement of transverse energy per unit pseudorapidity, \detdeta, in \auau collisions at \snn = 200 GeV using the sPHENIX calorimeter system.
Measurements of \detdeta are made by using layers of the sPHENIX calorimeter system separately, the Electromagnetic Calorimeter (EMCal), the Inner Hadronic Calorimeter (Inner HCal), and the Outer Hadronic Calorimeter (Outer HCal),
and by combining the energy deposits in all three calorimeter layers to make a full calorimeter measurement. 
Exploiting the nearly five hadronic interaction lengths of the full sPHENIX calorimeter system, this full calorimeter measurement can offer enhanced sensitivity to the total energy in \auau collisions compared to previous measurements that relied solely on electromagnetic energy. 

The \detdeta measurements presented here cover a pseudorapidity range of $|\eta| < 1.1$ for a range of Au+Au centrality intervals from $0\mathrm{-}70\%$ centrality. 
The average \detdeta as a function of the number of participating nucleons, \Npart, is also extracted for comparison to a variety of Monte Carlo heavy-ion event generators.
These measurements are the first physics results using the sPHENIX calorimeter system, and represent an important milestone in the commissioning and calibration of the detector.
Finally, the \detdeta results are in agreement with previous measurements at RHIC, and feature an improved granularity in $\eta$ and improved precision in low-$\Npart$ events.

The datasets used in this analysis along with reconstruction details are described in Section~\ref{sec:detdeta_datasets}.
The event selection criteria and centrality determination are described in Section~\ref{sec:detdeta_selection}.
The Monte Carlo (MC) simulation samples used in this analysis for correcting the measured \detdeta distributions back to the truth particle level are outlined in Section~\ref{sec:detdeta_mcsamples}.
The analysis procedure and determination of \detdeta correction factors are outlined in Section~\ref{sec:detdeta_anacorr}, and the systematic uncertainties associated with the measurement are described in Section~\ref{sec:detdeta_systematics}.
Finally, the results of the \detdeta measurements in \auau collisions at \snn = 200 GeV are presented in Section~\ref{sec:detdeta_results}.

% datasets and reconstruction -> Data and MC samples 
% event selections, trigger and reconstruction efficiency, background subtraction, unfolding, systematic uncertainties, final results
% jet reconstruction, data selection, trigger selection, monte carlo simulation, results, systematic uncertainties

% datasets and reconstruction 
% event selection
    % event selection
    % event centrality determination
% MC samples 
    % MC particle composition reweighting
    % MC z-vertex reweighting
% analysis procedure and correction factor determination
    % calorimeter tower energy reconstruction <- nominal method described in Chapter 3
    % vertex dependent eta correction
    % correction factors 
% systematics 
% results 

